\documentclass[11pt,a4paper]{article}

% -----------------------------------------------------------------------------
% Rapport de TP GitOps avec Argo CD
% Remplacer les champs entre crochets et les cadres de capture par vos éléments.
% Compilation : pdflatex report.tex (deux fois si nécessaire)
% -----------------------------------------------------------------------------
\usepackage[utf8]{inputenc}
\usepackage[T1]{fontenc}
\usepackage[french]{babel}
\usepackage{lmodern}
\usepackage[a4paper,margin=2.3cm]{geometry}
\usepackage{graphicx}
\usepackage{xcolor}
\usepackage{float}
\usepackage{booktabs}
\usepackage{tabularx}
\usepackage{array}
\usepackage{listings}
\usepackage{enumitem}
\usepackage{hyperref}
\usepackage{fancyhdr}
\usepackage{titlesec}

\definecolor{codegray}{HTML}{F4F4F4}
\definecolor{codegreen}{HTML}{2E7D32}
\definecolor{codeblue}{HTML}{1565C0}
\definecolor{codepurple}{HTML}{6A1B9A}
\definecolor{argored}{HTML}{C62828}

\lstdefinestyle{terminal}{
  backgroundcolor=\color{codegray},
  basicstyle=\ttfamily\small,
  breaklines=true,
  frame=single,
  rulecolor=\color{gray!40},
  columns=fullflexible,
  showstringspaces=false,
  commentstyle=\color{codegreen},
  keywordstyle=\color{codeblue},
  stringstyle=\color{codepurple}
}
\lstset{style=terminal}

\hypersetup{
  colorlinks=true,
  linkcolor=codeblue,
  urlcolor=codeblue,
  citecolor=codeblue,
  pdftitle={Mise en œuvre du Continuous Deployment avec GitOps et Argo CD},
  pdfauthor={[Nom de l'étudiant]}
}

\pagestyle{fancy}
\fancyhf{}
\lhead{TP GitOps et Argo CD}
\rhead{[Nom de l'étudiant]}
\cfoot{\thepage}
\setlength{\headheight}{14pt}
\titleformat{\section}{\Large\bfseries\color{argored}}{\thesection}{1em}{}
\titleformat{\subsection}{\large\bfseries\color{codeblue}}{\thesubsection}{1em}{}

\newcommand{\placeholder}[2][0.72\textwidth]{%
  \begin{figure}[H]
    \centering
    \fcolorbox{gray}{gray!8}{%
      \begin{minipage}[c][#1][c]{0.90\textwidth}
        \centering\small\textcolor{gray}{\textbf{PLACEHOLDER DE CAPTURE D'ÉCRAN}}\\[0.5em]
        #2\\[0.5em]
        \textit{Insérer ici la capture correspondante.}
      \end{minipage}%
    }
    \caption{#2}
  \end{figure}%
}

\newcommand{\file}[1]{\texttt{#1}}
\newcommand{\code}[1]{\texttt{#1}}

\begin{document}

\begin{titlepage}
  \centering
  \vspace*{1.5cm}
  {\Large \textsc{[Nom de l'établissement]}\par}
  \vspace{1.2cm}
  {\large Module DevOps / Kubernetes\par}
  \vspace{2cm}
  {\Huge\bfseries Mise en œuvre du Continuous Deployment\par}
  \vspace{0.35cm}
  {\Huge\bfseries avec GitOps et Argo CD\par}
  \vspace{1.5cm}
  {\Large Rapport de TP\par}
  \vfill
  \begin{tabular}{rl}
    Étudiant(e) : & [Nom et prénom] \\
    Classe : & [Classe / groupe] \\
    Enseignant(e) : & [Nom de l'enseignant(e)] \\
    Date : & 16 octobre 2026
  \end{tabular}
  \vfill
  \placeholder[3.5cm]{Vue générale de l'application Argo CD}
\end{titlepage}

\tableofcontents
\newpage

\section{Introduction}

Ce TP a pour objectif de mettre en œuvre un pipeline de déploiement continu suivant les principes GitOps. L'application choisie est un serveur NGINX simple déployé sur Kubernetes. Le dépôt Git contient la description déclarative de l'état souhaité et Argo CD surveille ce dépôt afin de réconcilier automatiquement cet état avec le cluster.

L'architecture réalisée est la suivante :

\begin{center}
\fbox{\parbox{0.86\textwidth}{
\centering
Commit dans Git $\longrightarrow$ Argo CD (pull) $\longrightarrow$ Cluster Kubernetes\\
\vspace{0.3em}
\small Argo CD compare l'état désiré et l'état réel, puis applique les changements nécessaires.
}}
\end{center}

Les objectifs sont de démontrer :
\begin{itemize}[itemsep=0.2em]
  \item le principe de \textit{Single Source of Truth} ;
  \item le déploiement automatique basé sur Git, sans push direct vers Kubernetes ;
  \item la détection du drift et le mécanisme de self-healing ;
  \item la différence entre le statut de synchronisation et le statut de santé ;
  \item l'historique des révisions et le rollback vers une version stable.
\end{itemize}

\section{Environnement technique}

\begin{table}[H]
\centering
\begin{tabularx}{\textwidth}{>{\bfseries}l X}
\toprule
Composant & Description \\
\midrule
Application & NGINX, image initiale \code{nginx:1.27-alpine} \\
Orchestrateur & Kubernetes [version] \\
CD / GitOps & Argo CD [version] \\
Packaging & Kustomize avec une base et un overlay \code{dev} \\
Namespace applicatif & \code{gitops-demo} \\
Dépôt Git & \url{https://github.com/youssef-ar/devops_lab1} \\
Cluster & [kind / Minikube / k3d / autre] \\
Système & [Linux / autre] \\
\bottomrule
\end{tabularx}
\caption{Environnement utilisé pour le TP.}
\end{table}

\subsection{Prérequis et installation}

Les outils suivants ont été utilisés : \code{kubectl}, \code{git}, Kustomize et la CLI Argo CD. Argo CD a été installé dans le namespace \code{argocd} avec le manifeste officiel :

\begin{lstlisting}[language=bash]
kubectl create namespace argocd
kubectl apply -n argocd \\
  -f https://raw.githubusercontent.com/argoproj/argo-cd/stable/manifests/install.yaml
kubectl wait --for=condition=Available deployment/argocd-server \\
  -n argocd --timeout=180s
\end{lstlisting}

L'interface a été rendue accessible localement avec un port-forward :

\begin{lstlisting}[language=bash]
kubectl -n argocd port-forward svc/argocd-server 8080:443
\end{lstlisting}

\placeholder{Installation d'Argo CD et Pods du namespace \code{argocd}}

\section{Structure du dépôt et configuration}

Le dépôt est organisé autour de Kustomize. Cette structure sépare les ressources communes de la configuration de l'environnement :

\begin{lstlisting}
.
├── apps/nginx/base/
│   ├── deployment.yaml
│   ├── service.yaml
│   └── kustomization.yaml
├── apps/nginx/overlays/dev/
│   └── kustomization.yaml
├── argocd/
│   ├── application.yaml
│   └── namespace.yaml
├── scripts/render.sh
└── README.md
\end{lstlisting}

Le Deployment décrit trois replicas dans l'état testé, ainsi que les probes de readiness et de liveness. Le Service expose NGINX dans le cluster sur le port 80. L'overlay \code{dev} ajoute le namespace \code{gitops-demo}.

\subsection{Application Argo CD}

L'objet Argo CD \code{nginx-gitops} pointe vers le dépôt suivant :

\begin{lstlisting}[language=yaml]
source:
  repoURL: https://github.com/youssef-ar/devops_lab1
  targetRevision: main
  path: apps/nginx/overlays/dev

destination:
  server: https://kubernetes.default.svc
  namespace: gitops-demo
\end{lstlisting}

La synchronisation automatique et le self-healing sont activés :

\begin{lstlisting}[language=yaml]
syncPolicy:
  automated:
    prune: true
    selfHeal: true
    allowEmpty: false
  syncOptions:
    - CreateNamespace=true
    - PruneLast=true
\end{lstlisting}

\placeholder{Application \code{nginx-gitops} affichée dans l'interface Argo CD}

\section{Déploiement initial}

Après publication du dépôt, l'application Argo CD a été créée avec :

\begin{lstlisting}[language=bash]
kubectl apply -f argocd/namespace.yaml
kubectl apply -f argocd/application.yaml

argocd app get nginx-gitops
kubectl get deploy,pods,svc -n gitops-demo
\end{lstlisting}

L'état attendu après le déploiement initial est :

\begin{itemize}
  \item application Argo CD : \code{Synced} ;
  \item santé de l'application : \code{Healthy} ;
  \item Deployment : trois replicas disponibles ;
  \item Pods : trois Pods \code{1/1 Running} ;
  \item Service : \code{ClusterIP} nommé \code{nginx}.
\end{itemize}

\placeholder{État initial : application Synced et Healthy dans Argo CD}

\section{Scénario 1 : Continuous Deployment}

\subsection{Modification de l'état désiré}

Pour tester le déploiement continu, le nombre de replicas a été modifié dans \file{apps/nginx/base/deployment.yaml}, puis le changement a été envoyé dans Git :

\begin{lstlisting}[language=bash]
sed -i 's/replicas: 2/replicas: 3/' \\
  apps/nginx/base/deployment.yaml

git add apps/nginx/base/deployment.yaml
git commit -m "feat: scale nginx to three replicas"
git push
\end{lstlisting}

Argo CD a détecté le nouveau commit et a appliqué le changement automatiquement. Aucune commande \code{kubectl apply} ni \code{argocd app sync} n'a été exécutée pour effectuer le déploiement.

\subsection{Résultat observé}

La commande suivante permet de suivre le déploiement :

\begin{lstlisting}[language=bash]
kubectl get deployment nginx -n gitops-demo -w
\end{lstlisting}

Le passage temporaire à \code{OutOfSync}, suivi du retour à \code{Synced}, confirme que le changement Git a été pris en compte. Le nombre de Pods disponibles atteint ensuite trois.

\placeholder{Argo CD détecte le commit et passe l'application à \code{Synced}}

\placeholder{Terminal montrant trois Pods NGINX disponibles après le déploiement}

\section{Scénario 2 : Drift et self-healing}

\subsection{Modification manuelle du cluster}

Afin de simuler une modification non autorisée, le Deployment a été supprimé directement du cluster :

\begin{lstlisting}[language=bash]
kubectl delete deployment nginx -n gitops-demo
kubectl get deployment nginx -n gitops-demo -w
\end{lstlisting}

Cette modification n'a pas été effectuée dans Git. Elle représente donc un écart entre l'état réel du cluster et l'état désiré décrit dans le dépôt.

\subsection{Réconciliation automatique}

Grâce à \code{selfHeal: true}, Argo CD a détecté que le Deployment était absent et l'a recréé automatiquement. Le dépôt Git n'a pas été modifié : il est resté la source de vérité.

\placeholder{Argo CD affiche temporairement le drift : ressource Missing ou OutOfSync}

\placeholder{Recréation automatique du Deployment par Argo CD}

\section{Scénario 3 : Sync status et Health status}

\subsection{Introduction d'une erreur applicative}

Une image inexistante a été introduite volontairement dans le manifeste :

\begin{lstlisting}[language=yaml]
image: nginx:does-not-exist
\end{lstlisting}

Après commit et push, Argo CD a correctement appliqué la déclaration provenant de Git. Cependant, Kubernetes n'a pas pu démarrer le conteneur, car l'image n'existe pas dans le registre :

\begin{lstlisting}[language=bash]
kubectl get pods -n gitops-demo
\end{lstlisting}

Résultat observé :

\begin{lstlisting}
NAME                     READY   STATUS         RESTARTS   AGE
nginx-5994554846-vx6wv   0/1     ErrImagePull   0          67s
nginx-c4b58fbc-6gd4h     1/1     Running        0          2m48s
nginx-c4b58fbc-bmrb4     1/1     Running        0          2m48s
nginx-c4b58fbc-nm62b     1/1     Running        0          2m48s
\end{lstlisting}

Les anciens Pods fonctionnels sont conservés pendant le rolling update afin de limiter l'interruption de service. Le nouveau ReplicaSet ne devient pas disponible, car son Pod est en \code{ErrImagePull}.

\subsection{Interprétation des statuts}

\begin{table}[H]
\centering
\begin{tabularx}{\textwidth}{>{\bfseries}l X}
\toprule
Statut & Interprétation \\
\midrule
Sync = Synced & Les manifests présents dans le cluster correspondent à l'état défini dans Git. \\
Health = Degraded / Progressing & L'application déclarée est synchronisée, mais son fonctionnement réel est dégradé. \\
Pod = ErrImagePull & Le kubelet ne parvient pas à télécharger l'image du conteneur. \\
\bottomrule
\end{tabularx}
\caption{Différence entre synchronisation et santé applicative.}
\end{table}

Le diagnostic détaillé a été obtenu avec :

\begin{lstlisting}[language=bash]
kubectl describe pod nginx-5994554846-vx6wv -n gitops-demo
kubectl get events -n gitops-demo --sort-by=.lastTimestamp
\end{lstlisting}

\placeholder{Argo CD : application Synced mais état Health Degraded ou Progressing}

\placeholder{Terminal : Pod en ErrImagePull ou ImagePullBackOff}

\section{Scénario 4 : Historique et rollback}

Depuis l'interface Argo CD, l'historique de l'application \code{nginx-gitops} permet de consulter les révisions associées aux commits Git. La révision stable précédente a été sélectionnée dans la page \textit{History and Rollback}, puis un rollback a été lancé.

\begin{enumerate}
  \item Ouvrir l'application \code{nginx-gitops} dans Argo CD ;
  \item ouvrir la vue \textit{History and Rollback} ;
  \item sélectionner la dernière révision stable ;
  \item lancer \textit{Rollback} ;
  \item vérifier le retour des Pods à l'état \code{Running}.
\end{enumerate}

Comme l'auto-synchronisation est activée, le rollback d'Argo CD doit être considéré comme une opération de retour rapide. La correction durable consiste à remettre une image valide dans Git :

\begin{lstlisting}[language=bash]
sed -i 's#nginx:does-not-exist#nginx:1.27-alpine#' \\
  apps/nginx/base/deployment.yaml

git add apps/nginx/base/deployment.yaml
git commit -m "fix: restore stable nginx image"
git push
\end{lstlisting}

Cette correction permet de réaligner durablement le dépôt et le cluster.

\placeholder{Vue \textit{History and Rollback} avec les révisions de l'application}

\placeholder{Rollback vers la révision stable précédente}

\placeholder{Après rollback : application Healthy et Pods Running}

\section{Bilan et analyse}

Les objectifs du TP sont couverts comme suit :

\begin{table}[H]
\centering
\begin{tabularx}{\textwidth}{>{\bfseries}p{0.27\textwidth} X}
\toprule
Concept & Mise en œuvre et preuve \\
\midrule
Single Source of Truth & Les ressources Kubernetes sont versionnées dans le dépôt Git et l'Application Argo CD pointe vers ce dépôt. \\
Pull-based deployment & Argo CD surveille le dépôt et applique les commits ; aucun déploiement direct depuis l'extérieur du cluster n'est nécessaire. \\
Configuration drift & La suppression manuelle du Deployment crée un écart détecté par Argo CD. \\
Self-healing & L'option \code{selfHeal: true} recrée automatiquement la ressource absente. \\
Sync vs Health & Une image invalide produit une déclaration \code{Synced}, mais un Pod en \code{ErrImagePull} et une santé dégradée. \\
Rollback & L'historique Argo CD permet de revenir rapidement à une révision stable ; la correction finale est commitée dans Git. \\
\bottomrule
\end{tabularx}
\caption{Correspondance entre les concepts demandés et les résultats du TP.}
\end{table}

Ce TP montre que GitOps sépare clairement la déclaration de l'état souhaité, stockée dans Git, de la réconciliation réalisée par Argo CD. Cette approche améliore la traçabilité, la reproductibilité et la capacité à restaurer rapidement un état stable.

\section{Conclusion}

La mise en place de l'application NGINX a permis de démontrer un pipeline GitOps complet avec Kustomize, Kubernetes et Argo CD. Le dépôt Git joue le rôle de source unique de vérité et Argo CD assure la synchronisation continue avec le cluster. Les tests de modification Git, de drift, d'image invalide et de rollback illustrent les mécanismes essentiels du Continuous Deployment pull-based.

La principale limite observée est qu'un état peut être techniquement synchronisé avec Git tout en étant incorrect du point de vue applicatif. Le contrôle du statut de santé et l'ajout de tests de validation dans une chaîne CI sont donc nécessaires pour compléter cette architecture dans un contexte de production.

\section*{Annexe : commandes utiles}
\addcontentsline{toc}{section}{Annexe : commandes utiles}

\begin{lstlisting}[language=bash]
# Vérification locale du rendu Kustomize
kubectl kustomize apps/nginx/overlays/dev
kubectl apply --dry-run=client -k apps/nginx/overlays/dev

# État Kubernetes
kubectl get all -n gitops-demo
kubectl get events -n gitops-demo --sort-by=.lastTimestamp

# État et historique Argo CD
argocd app get nginx-gitops
argocd app history nginx-gitops
\end{lstlisting}

\end{document}
